% ch07_u11 —— 原书 p.336–p.345（第 7 章 7.13 BRST 基础）
% 底本：A. Beilinson, V. Drinfeld, Quantization of Hitchin's Integrable System
% and Hecke Eigensheaves

\SourcePage{336}

\emph{注.} 在 (d) 中我们不得不排除 $A$ 为完备局部环的情形。真正的原因由下述结果说明：
\begin{enumerate}[label=\arabic*),leftmargin=2.6em]
\item 据 [J]，若 $A$ 是完备局部 Noether 环，$M$ 是平坦 $A$-模，$N$ 是有限生成 $A$-模，则 $\operatorname{Ext}(M,N)=0$；
\item 据 [RG]（p.76，2.3.3 的注 4），若 $A$ 是 Artin 环的投射极限（这是 [RG] 中「线性紧」一词的含义吗？），则每个（平坦？）Mittag-Leffler $A$-模都是严格 Mittag-Leffler 的。（[RG] 中没有平坦性假设，但去掉这一假设后他们的论证还正确吗？例如，[RG] 中的 $F_i$ 为何是线性紧的？）
\end{enumerate}

\bdsec{7.13. BRST 基础}

BRST 构造是哈密顿约化的精致版本，它在无穷维情形尤为重要。本文正文两次调用 BRST：先是用它定义 Feigin--Frenkel 同构，随后用它构造 Hecke 性质证明中所用的局部化函子 $L\Delta$。本节简要叙述一般的 BRST 构造；函子 $L\Delta$ 将在下一节研究。

BRST 的常用数学参考文献是 [F84]、[FGZ86]、[KS] 与 [Ak]。我们试图写一份不依赖冗余结构的叙述（如 $\mathbb{Z}$-分次、正规序等）。
% [Tutor Note: 原书 p.336 此处印作 "redundand structures"，"redundand" 系 "redundant" 之误，底本如此，保留原样。]

我们先从有限维情形开始。在关于 Tate 中心扩张的一段插叙之后，再解释无穷维版本。

\bdp{7.13.1} 设 $F$ 是有限维向量空间。以 $\mathrm{Cl}^{\cdot}=\mathrm{Cl}^{\cdot}_F$ 记 $F\oplus F^{*}$ 的 Clifford 代数，其分次使 $F$ 的度为 $-1$、$F^{*}$ 的度为 $1$。我们把 $\mathrm{Cl}^{\cdot}$ 视为分次向量空间的张量范畴中的代数\footnote{带「超」交换约束。}。置 $\mathrm{Cl}^{\cdot}_i:=\Lambda^{\le i}F\cdot\Lambda F^{*}\subset\mathrm{Cl}^{\cdot}$。则 $\mathrm{Cl}^{\cdot}_0=\Lambda^{\cdot}F^{*}\subset\mathrm{Cl}^{\cdot}_1\subset\ldots$ 是 $\mathrm{Cl}^{\cdot}$ 上的环滤链。\emph{经典 Clifford 代数} $\mathcal{C}l^{\cdot}=\mathcal{C}l^{\cdot}_F:=\operatorname{gr}\mathrm{Cl}^{\cdot}$ 是交换的（作为分次代数），故它是\SourcePage{337}%
通常意义下的泊松代数。

置 $\mathcal{C}l^{\cdot}_i:=\operatorname{gr}_i\mathrm{Cl}^{\cdot}$。交换分次代数 $\mathcal{C}l^{\cdot}$ 由 $F=\mathcal{C}l^{-1}_1$ 与 $F^{*}=\mathcal{C}l^{0}_1$ 自由生成。泊松括号 $\{\,,\,\}$ 在 $F$ 与 $F^{*}$ 上为零，而对 $f\in F$、$f^{*}\in F^{*}$ 有 $\{f,f^{*}\}=f^{*}(f)$。

子空间 $\mathcal{C}l^{0}_1$ 是 $\mathcal{C}l^{\cdot}$ 的李子代数；它使 $F$ 与 $F^{*}$ 正规化，相应的伴随作用把它与 $\operatorname{End}_F$ 及 $\operatorname{End}_{F^{*}}$ 等同。令 $E^{Lie}=\operatorname{End}^{Lie}_F$ 表示把 $\operatorname{End}_F$ 视为李代数。则 $E^{\flat}=\operatorname{End}^{\flat}_F:=\mathrm{Cl}^{0}_1$ 是 $E^{Lie}$ 被 $\mathbb{C}$ 的中心扩张。

\emph{注.} (i) $\mathrm{Cl}$ 在 $\Lambda F^{*}\xrightarrow{\sim}\mathrm{Cl}/\mathrm{Cl}\cdot F$ 上的作用把它与「奇」向量空间 $F^{odd}$ 上的微分算子代数等同。$\mathrm{Cl}$ 上的滤链就是通常按微分算子阶数给出的滤链，故 $\mathcal{C}l^{\cdot}$ 是 $F^{odd}$ 的余切丛上函数的泊松代数。

(ii)（只在有限维情形有效）扩张 $\operatorname{End}^{\flat}_F$ 分裂（分裂方式不唯一）。事实上，有分裂 $s',s'':E^{Lie}\to E^{\flat}$，它们分别把 $E^{Lie}$ 与 $F^{*}\cdot F$ 及 $F\cdot F^{*}$ 等同。其余任一分裂都等于 $s_{\lambda}=\lambda s'+(1-\lambda)s''$，其中 $\lambda\in\mathbb{C}$。例如 $s_{1/2}$ 是「酉」分裂，它也可如下定义。注意 $\mathrm{Cl}$ 带有典范反自同构（作为\emph{分次}代数），它在 $F$ 与 $F^{*}$ 上是恒等。它保持 $\mathrm{Cl}^{0}_1$，而「酉」分裂就是 $-1$ 特征子空间。

\bdp{7.13.2} 这里是 BRST 构造的「经典」版本。设 $\mathfrak{n}$ 是有限维李代数，$\mathcal{R}$ 是泊松代数，$l^{c}:\mathfrak{n}\to\mathcal{R}$ 是李代数态射\footnote{「c」表示「classical」（经典）。}。置 $\mathcal{C}l^{\cdot}:=\mathcal{C}l^{\cdot}_{\mathfrak{n}}$。$\mathfrak{n}$ 的伴随作用给出李代数态射 $a^{c}:\mathfrak{n}\to\mathcal{C}l^{0}_1$。置 $\mathcal{A}^{\cdot}:=\mathcal{C}l^{\cdot}\otimes\mathcal{R}$；这是泊松分次代数。它还带有附加分次 $\mathcal{A}^{\cdot}_{(i)}:=\mathcal{C}l^{\cdot}_i\otimes\mathcal{R}$，它与乘法相容（但与泊松括号不相容）。我们有李代数态射 $\mathcal{L}ie:\mathfrak{n}\to\mathcal{A}^{0}$，$n\mapsto\mathcal{L}ie_n:=a^{c}(n)\otimes 1+1\otimes l^{c}(n)$。下文对 $n\in\mathfrak{n}$，以 $i^{c}_n$ 记 $\mathcal{C}l^{-1}_1\subset\mathcal{A}^{-1}_{(1)}$ 中相应的元素。有 $\{\mathcal{L}ie_{n_1},i^{c}_{n_2}\}=i^{c}_{[n_1,n_2]}$。

\SourcePage{338}

下列关键引理及其「量子」版本 7.13.7 本质上属于 Akman [Ak]。

\bdp{7.13.3} \emph{引理.} 存在唯一元素 $Q^{c}=Q^{c}_{\mathcal{A}}\in\mathcal{A}^{1}$，使得对任意 $n\in\mathfrak{n}$ 有 $\{Q^{c},i^{c}_n\}=\mathcal{L}ie_n$。事实上 $Q^{c}\in\mathcal{A}^{1}_{(\le1)}$。有 $\{Q^{c},Q^{c}\}=0$。

\emph{证明.} 把 $\mathcal{A}$ 视为 $\Lambda\mathfrak{n}$-模，其中 $n\in\mathfrak{n}=\Lambda^{1}\mathfrak{n}$ 作用为 $\operatorname{Ad}_{i^{c}_n}=\{i^{c}_n,\cdot\}$。被所有 $\operatorname{Ad}_{i^{c}_n}$ 零化的元素构成的子空间（即 $\mathfrak{n}\subset\mathcal{A}^{-1}_{(1)}$ 的中心化子）等于 $\Lambda\mathfrak{n}\otimes\mathcal{R}$。它是 $\mathcal{A}^{\le0}$ 的子空间，故 $Q^{c}$ 的唯一性显然。我们的 $\Lambda\mathfrak{n}$-模是自由的，故 $Q^{c}$ 的存在性由下述事实推出：映射 $n_1,n_2\mapsto\{\mathcal{L}ie_{n_1},i^{c}_{n_2}\}$ 是斜对称的。因为 $\mathcal{L}ie_n\in\mathcal{A}^{0}_{(\le1)}$，所以 $Q^{c}\in\mathcal{A}^{1}_{(\le1)}$。最后，因 $\{Q^{c},Q^{c}\}\in\mathcal{A}^{2}$，要验证它为零，只需证明对任意 $n,n'\in\mathfrak{n}$ 有 $\operatorname{Ad}_{i^{c}_n}\operatorname{Ad}_{i^{c}_{n'}}(\{Q^{c},Q^{c}\})=0$。事实上，
\[
\operatorname{Ad}_{i^{c}_n}\operatorname{Ad}_{i^{c}_{n'}}(\{Q^{c},Q^{c}\})
=2\operatorname{Ad}_{i^{c}_n}(\{\mathcal{L}ie_{n'},Q^{c}\})
=2\{i^{c}_{[n,n']},Q^{c}\}+2\{\mathcal{L}ie_{n'},\mathcal{L}ie_n\}=0.
\]
\hfill$\Box$

\emph{注.} 以 $\mathfrak{n}_{\heartsuit}$ 记下述李分次代数：其非零分量为 $\mathfrak{n}^{-1}_{\heartsuit}=\mathfrak{n}$、$\mathfrak{n}^{0}_{\heartsuit}=\mathfrak{n}$、$\mathfrak{n}^{1}_{\heartsuit}=\mathbb{C}=\mathbb{C}\cdot Q$；$\mathfrak{n}^{0}_{\heartsuit}$ 上的李括号与 $\mathfrak{n}$ 的相同；$\mathfrak{n}^{0}_{\heartsuit}$ 在 $\mathfrak{n}^{-1}_{\heartsuit}$ 上的伴随作用是 $\mathfrak{n}$ 的伴随作用；算子 $\operatorname{Ad}_{Q}:\mathfrak{n}^{-1}_{\heartsuit}\to\mathfrak{n}^{0}_{\heartsuit}$ 是 $\operatorname{id}_{\mathfrak{n}}$。于是带微分 $\operatorname{Ad}_{Q}$ 的 $\mathfrak{n}_{\heartsuit}$ 是李 DG 代数\footnote{注意 $\mathfrak{n}_{\heartsuit}/\mathfrak{n}^{1}_{\heartsuit}$ 就是 7.6.3 中的李 DG 代数 $\mathfrak{n}_{\Omega}$。}。此时 7.13.3 说明：存在李分次代数的典范态射 $\mathcal{L}ie:\mathfrak{n}_{\heartsuit}\to\mathcal{A}^{\cdot}$，其分量分别为 $n\mapsto i^{c}_n$、$n\mapsto\mathcal{L}ie_n$、$Q\mapsto Q^{c}$。

\bdp{7.13.4} 置 $d:=\operatorname{Ad}_{Q^{c}}=\{Q^{c},\cdot\}$。这是 $\mathcal{A}^{\cdot}$ 的度 $1$ 的导子，且平方为零。于是 $\mathcal{A}$ 是泊松 DG 代数；它称为 $\mathcal{R}$ 的\emph{BRST 约化}。态射 $\mathcal{L}ie:\mathfrak{n}_{\heartsuit}\to\mathcal{A}$ 是李 DG 代数态射。

若 $H^{i}\mathcal{A}=0$（$i\ne0$），则称 BRST 约化是\emph{正则的}。

易见 $Q^{c}=Q_1+Q_0$，其中 $Q_1\in\mathcal{A}^{1}_{(1)}=\mathfrak{n}\otimes\Lambda^{2}\mathfrak{n}^{*}\otimes\mathcal{R}$ 与 $Q_0\in\mathcal{A}^{1}_{(0)}=\mathfrak{n}^{*}\otimes\mathcal{R}$ 分别是 $\frac{1}{2}a^{c}\in\operatorname{Hom}(\mathfrak{n},\mathcal{C}l^{0}_1)=\mathfrak{n}^{*}\otimes\mathcal{C}l^{0}_1\subset\mathcal{A}^{1}\otimes\mathcal{A}^{0}$ 在乘法映射下的像，而 $l\in\operatorname{Hom}(\mathfrak{n},\mathcal{R})=\mathcal{A}^{1}_{(1)}$。
% [Tutor Note: 原书 p.338 此处印作 "$l\in\operatorname{Hom}(\mathfrak{n},\mathcal{R})=\mathcal{A}^{1}_{(1)}$"；因 $\operatorname{Hom}(\mathfrak{n},\mathcal{R})=\mathfrak{n}^{*}\otimes\mathcal{R}=\mathcal{A}^{1}_{(0)}$，末项 $(1)$ 疑为 $(0)$ 之误，底本如此，保留原样。]
把微分按双分次分解，\SourcePage{339}%
可见 $\mathcal{A}$ 是以 $d':\mathcal{A}^{i}_{(j)}\to\mathcal{A}^{i+1}_{(j)}$、$d'':\mathcal{A}^{i}_{(j)}\to\mathcal{A}^{i+1}_{(j-1)}$ 为双微分的双复形的全复形。

BRST 微分保持滤链 $\mathcal{A}_{(\le i)}$。特别地 $\mathcal{A}_{(0)}=C(\mathfrak{n},\mathcal{R})$ 是 $\mathcal{A}$ 的 DG 子代数，于是有分次代数的典范态射
\begin{equation}\tag{365}
H^{\cdot}(\mathfrak{n},\mathcal{R})\to H^{\cdot}\mathcal{A}.
\end{equation}

注意 $(\mathcal{A}^{\cdot}_{(-)},d'')$ 是关于 $l^{c}:\mathfrak{n}\to\mathcal{R}$ 的 Koszul 复形 $P:=\Lambda^{-\cdot}\mathfrak{n}\otimes\mathcal{R}$。于是 $\mathcal{A}$ 是 $\mathfrak{n}$ 的以 $P$ 为系数的李代数上链的 Chevalley 复形 $C^{\cdot}(\mathfrak{n},P)$。显然的投影 $P\to\mathcal{R}/\mathcal{R}l^{c}(\mathfrak{n})$ 给出 DG 代数同构 $\mathcal{A}/\mathcal{I}\xrightarrow{\sim}C(\mathfrak{n},\mathcal{R}/\mathcal{R}l^{c}(\mathfrak{n}))$，其中 $\mathcal{I}\subset\mathcal{A}$ 是由元素 $i^{c}_n$（$n\in\mathfrak{n}$）生成的 DG 理想。过渡到上同调，得到分次代数的典范态射
\begin{equation}\tag{366}
H^{\cdot}\mathcal{A}\to H^{\cdot}(\mathfrak{n},\mathcal{R}/\mathcal{R}l^{c}(\mathfrak{n})).
\end{equation}

若 $H_i(P)=0$（$i\ne0$），则称 $l^{c}$ 是\emph{正则的}。

\bdp{7.13.5} \emph{引理.} 若 $l^{c}$ 正则，则 (366) 是同构。

\emph{证明.} 正则性意味着投影 $P\to\mathcal{R}/\mathcal{R}l^{c}(\mathfrak{n})$ 是拟同构。故 $\mathcal{A}\to C^{\cdot}(\mathfrak{n},\mathcal{R}/\mathcal{R}l^{c}(\mathfrak{n}))$ 也是拟同构。\hfill$\Box$

于是对负的 $i$ 有 $H^{i}\mathcal{A}=0$，且 $H^{0}\mathcal{A}\xrightarrow{\sim}[\mathcal{R}/\mathcal{R}l^{c}(\mathfrak{n})]^{\mathfrak{n}}$，这正是 $\mathcal{R}$ 关于哈密顿作用 $l^{c}$ 的通常哈密顿约化。

\bdp{7.13.6} 现在转向 BRST 的「量子」版本。设 $\mathfrak{n}$ 是有限维李代数。置 $\mathrm{Cl}^{\cdot}:=\mathrm{Cl}^{\cdot}_{\mathfrak{n}}$。以 $\mathfrak{n}^{\flat}$ 记 $\mathfrak{n}$ 被 $\mathbb{C}$ 的中心扩张，它定义为 $\operatorname{End}^{\flat}_{\mathfrak{n}}$ 沿伴随作用态射 $\mathfrak{n}\to\operatorname{End}_{\mathfrak{n}}$ 的拉回（记号见 7.13.1 末）。换言之，$\mathfrak{n}^{\flat}$ 是 $\mathfrak{n}$ 被 $\mathbb{C}$ 的中心扩张，配有李代数映射 $a:\mathfrak{n}^{\flat}\to\mathrm{Cl}^{0}$，使 $a(1_{\mathfrak{n}^{\flat}})=1$\footnote{这里 $1_{\mathfrak{n}^{\flat}}$ 是 $\mathbb{C}\subset\mathfrak{n}^{\flat}$ 的生成元。}，\SourcePage{340}%
且由 $\mathfrak{n}\oplus\mathfrak{n}^{*}$ 上的伴随作用在 $\mathrm{Cl}$ 上诱导的 $\mathfrak{n}$-作用与 $a$ 的伴随作用相同。

设 $R$ 是结合代数，$l:\mathfrak{n}^{\flat}\to R$ 是李代数态射，且 $l(1_{\mathfrak{n}^{\flat}})=-1$。置 $A^{\cdot}:=\mathrm{Cl}^{\cdot}\otimes R$；这是结合分次代数。我们有李代数态射 $\operatorname{Lie}:=a+l:\mathfrak{n}\to A^{0}$，$n\mapsto\operatorname{Lie}_n:=a(n^{\flat})+l(n^{\flat})$，其中 $n^{\flat}$ 是 $n$ 到 $\mathfrak{n}^{\flat}$ 的任一提升。下文对 $n\in\mathfrak{n}$，以 $i_n$ 记 $\mathrm{Cl}^{-1}_1\subset A^{-1}$ 中相应的元素。有 $[\operatorname{Lie}_{n_1},i_{n_2}]=i_{[n_1,n_2]}$。

\bdp{7.13.7} \emph{引理.} 存在唯一元素 $Q=Q_A\in A^{1}$，使得对任意 $n\in\mathfrak{n}$ 有 $[Q,i_n]=\operatorname{Lie}_n$。事实上 $Q\in\mathrm{Cl}^{1}_1\otimes R$。有 $Q^{2}=0$。

\emph{证明.} 与「经典」版本 7.13.3 的证明相同。\hfill$\Box$

置 $d:=\operatorname{Ad}_Q$\footnote{当然，我们取「超」意义下的 $\operatorname{Ad}$，故对 $v\in A^{odd}$ 有 $dv=Qv+vQ$。}；这是 $A$ 的度 $1$ 的导子，且平方为零。于是 $A$ 是结合 DG 代数，称为 $R$ 的\emph{BRST 约化}。如同 7.13.3 后的注与 7.13.4 那样，我们有李 DG 代数的典范态射 $\operatorname{Lie}:\mathfrak{n}_{\heartsuit}\to A$，其分量为 $n\mapsto i_n$、$n\mapsto\operatorname{Lie}_n$、$Q\mapsto Q_A$。

若 $H^{i}A=0$（$i\ne0$），则称 BRST 约化是\emph{正则的}。

以 $C(\mathfrak{n},R)$ 记 $\mathfrak{n}$ 的以 $R$ 为系数的李代数上链的 Chevalley DG 代数（关于作用 $\operatorname{Ad}_l$）。作为分次代数它等于 $\Lambda^{\cdot}\mathfrak{n}^{*}\otimes R$，故它是 $A^{\cdot}$ 的子代数。

\bdp{7.13.8} \emph{引理.} 嵌入 $C(\mathfrak{n},R)\subset A$ 与微分相容。

\emph{证明.} 只需说明：在 $R,\mathfrak{n}^{*}\subset A$ 上，我们的微分分别等于 $\mathfrak{n}$-作用映射 $R\to\mathfrak{n}^{*}\otimes R$ 的对偶与括号映射 $\mathfrak{n}^{*}\to\Lambda^{2}\mathfrak{n}^{*}$ 的对偶。如同 $Q$ 的唯一性证明，只需验证对任意 $n,n_1,n_2\in\mathfrak{n}$、$n^{*}\in\mathfrak{n}^{*}$、$r\in R$ 有 $[i_n,[Q,r]]=[l(n),r]$ 与 $[i_{n_1},[i_{n_2},[Q,n^{*}]]]=n^{*}([n_1,n_2])$；这是直接的计算。\hfill$\Box$
% [Tutor Note: 原书 p.340 此处印作 "It suuffices to show"，"suuffices" 系 "suffices" 之误，底本如此，保留原样。]

\emph{注.} 我们看到 $d$ 保持环滤链 $\mathrm{Cl}^{\cdot}\otimes R$。在 $\mathrm{Cl}_i\otimes R/\mathrm{Cl}_{i-1}\otimes R=\Lambda^{\cdot+i}\mathfrak{n}^{*}\otimes\Lambda^{\cdot}\mathfrak{n}\otimes R=C^{\cdot+i}(\mathfrak{n},\Lambda^{\cdot}\mathfrak{n}\otimes R)$ 上，它与 Chevalley 微分相同。

\SourcePage{341}

DG 代数嵌入 $C(\mathfrak{n},R)\subset A$ 给出分次代数态射
\begin{equation}\tag{367}
H^{\cdot}(\mathfrak{n},R)\to H^{\cdot}A.
\end{equation}
特别地，因 $R$ 的中心 $\mathfrak{z}$ 落在 $R^{\mathfrak{n}}$ 中，得到态射
\begin{equation}\tag{368}
\mathfrak{z}\to H^{0}A.
\end{equation}

\bdp{7.13.9} \emph{注.}（只在有限维情形有效）设 $I$ 是由元素 $i_n$（$n\in\mathfrak{n}$）生成的 $A$ 的左 DG 理想。商复形 $A/I$ 可如下计算。设 $\mathfrak{n}\hookrightarrow\mathfrak{n}^{\flat}$ 是由 7.13.1 的注 (ii) 中的分裂 $s'$ 定义的分裂。则 $I$ 作为普通理想由元素 $i_n$ 与 $l(n)$（$n\in\mathfrak{n}$）生成。把投影 $A\to A/I$ 限制到 $C(\mathfrak{n},R)$，得到复形的同构 $A/I\xrightarrow{\sim}C(\mathfrak{n},R/Rl(\mathfrak{n}))$，它给出态射
\begin{equation}\tag{369}
H^{\cdot}A\to H^{\cdot}(\mathfrak{n},R/Rl(\mathfrak{n})).
\end{equation}

\bdp{7.13.10} \emph{注.} 设 $C^{\cdot}$ 是不可约分次 $\mathrm{Cl}^{\cdot}$-模（这样的 $C^{\cdot}$ 在同构与分次平移的意义下唯一）。若 $M=(M^{\cdot},d_M)$ 是 $R$-复形（即 $R$-模的复形），则 $M\otimes C^{\cdot}:=(M^{\cdot}\otimes C^{\cdot},d)$，其中 $d:=d_M\otimes\operatorname{id}_{C^{\cdot}}+Q\cdot$，是 $A$-复形（即 DG $A$-模）。函子 $\cdot\otimes C^{\cdot}$：($R$-复形) $\to$ ($A$-复形) 是范畴等价。

\bdp{7.13.11} 我们来比较「量子」与「经典」两种情形。设处于 7.13.6 的情形。设 $R_0\subset R_1\subset\ldots$ 是 $R$ 上的递增环滤链，使 $\cup R_i=R$ 且 $\mathcal{R}:=\operatorname{gr}R$ 交换。则 $\mathcal{R}$ 按通常方式是泊松代数。给 $A$ 配上滤链 $A_{\cdot}$，它等于滤链 $\mathrm{Cl}^{\cdot}$ 与 $R_{\cdot}$ 的张量积。则 $\mathcal{A}:=\operatorname{gr}A$ 作为泊松分次代数等于 $\mathcal{C}l^{\cdot}\otimes\mathcal{R}$。置 $\mathcal{A}_i:=\operatorname{gr}_iA$。

设 $l(\mathfrak{n}^{\flat})\subset R_1$；以 $l^{c}$ 记相应的态射 $\mathfrak{n}\to\mathcal{R}_1$。则 $(\mathcal{R},l^{c})$ 是定义 7.13.2 中「经典」BRST 构造的数据。由 7.13.3 有相应的「经典」BRST 元素 $Q^{c}$。易见 $Q\in A_1$，且 $Q^{c}$ 等于 $Q$ 在 $\mathcal{A}_1$ 中的像。

\SourcePage{342}

因此滤链 $A_{\cdot}$ 关于微分稳定，且 $\operatorname{gr}A$ 作为泊松 DG 代数与相应的「经典」$\mathcal{A}$ 相同。于是有收敛到 $H^{\cdot}A$ 的谱序列，其首项为
\[
E^{p,q}_1=H^{p+q}\mathcal{A}_{-p}.
\]

\bdp{7.13.12} \emph{引理.} (i) 设 $l^{c}$ 正则。则对 $i<0$ 有 $H^{i}\mathcal{A}=0$，且 $\operatorname{gr}H^{0}\mathcal{A}\subset[\mathcal{R}/\mathcal{R}l^{c}(\mathfrak{n})]^{\mathfrak{n}}$。

(ii) 若此外还有 $H^{i}(\mathfrak{n},\mathcal{R}/\mathcal{R}l^{c}(\mathfrak{n}))=0$（$i>0$），则对 $i\ne0$ 有 $H^{i}\mathcal{A}=0$，且 $\operatorname{gr}H^{0}\mathcal{A}\xrightarrow{\sim}[\mathcal{R}/\mathcal{R}l^{c}(\mathfrak{n})]^{\mathfrak{n}}$。

\emph{证明.} 看谱序列与 7.13.5。\hfill$\Box$

\bdp{7.13.13} 在下述情形可以显式算出代数 $H^{0}A$。设处于 7.13.11 的情形，且 $l:\mathfrak{n}^{\flat}\to R_1$ 是单射。以 $\mathfrak{b}'$ 记 $l(\mathfrak{n}^{\flat})$ 在 $R_1$ 中的正规化子。于是 $\mathfrak{b}'$ 是含 $\mathfrak{n}^{\flat}$ 的李代数，且有延拓 $l$ 的李代数嵌入 $l^{\mathfrak{b}}:\mathfrak{b}'\to R_1$。置 $\mathfrak{b}:=\mathfrak{b}'/\mathbb{C}$，于是 $\mathfrak{b}'$ 是 $\mathfrak{b}$ 被 $\mathbb{C}$ 的中心扩张。$\mathfrak{b}$ 的伴随作用给出李代数态射 $\mathfrak{b}\to\operatorname{End}_{\mathfrak{n}}$；以 $\mathfrak{b}^{\flat}$ 记中心扩张 $\operatorname{End}^{\flat}$ 的拉回（见 7.13.1）。则 $\mathfrak{n}^{\flat}$ 是 $\mathfrak{b}^{\flat}$ 的李子代数，且有延拓 $a$ 的李代数态射 $a^{\mathfrak{b}}:\mathfrak{b}^{\flat}\to\mathrm{Cl}^{0}_1$。

设 $\mathfrak{b}^{\natural}$ 是扩张 $\mathfrak{b}'$ 与 $\mathfrak{b}^{\flat}$ 的 Baer 和。按构造有典范分裂 $s:\mathfrak{n}\to\mathfrak{b}^{\natural}$。它关于 $\mathfrak{b}$ 的伴随作用不变，故 $s(\mathfrak{n})$ 是 $\mathfrak{b}^{\natural}$ 中的理想。置 $\mathfrak{h}^{\natural}:=\mathfrak{b}^{\natural}/s(\mathfrak{n})$；这是 $\mathfrak{h}:=\mathfrak{b}/\mathfrak{n}$ 被 $\mathbb{C}$ 的中心扩张。

置 $\operatorname{Lie}^{\mathfrak{b}}:=a^{\mathfrak{b}}\otimes1+1\otimes l^{\mathfrak{b}}:\mathfrak{b}^{\natural}\to A^{0}_1$。这是李代数态射，它在 $\mathbb{C}\subset\mathfrak{b}^{\natural}$ 上等于 $\operatorname{id}_{\mathbb{C}}$。其像与 $Q$ 交换（因为我们的构造都是自然的），即它属于 $\operatorname{Ker}d$。有 $\operatorname{Lie}^{\mathfrak{b}}\circ s=\operatorname{Lie}=d\circ i:\mathfrak{n}\to A^{0}$，故 $\operatorname{Lie}^{\mathfrak{b}}$ 给出典范态射 $\operatorname{Lie}^{\mathfrak{h}}:\mathfrak{h}^{\natural}\to H^{0}A$。以 $U^{\natural}\mathfrak{h}$ 记 $\mathfrak{h}$ 的对应于 $\mathfrak{h}^{\natural}$ 的扭包络代数。我们的 $\operatorname{Lie}^{\mathfrak{h}}$ 给出结合代数的典范态射\SourcePage{343}%
\begin{equation}\tag{370}
h:U^{\natural}\mathfrak{h}\to H^{0}A.
\end{equation}

此态射有显然的「经典」版本 $h^{c}:\operatorname{Sym}\mathfrak{h}\to H^{0}\mathcal{A}$。它与投影 $H^{0}A\to[\mathcal{R}/l^{c}(\mathfrak{n})\mathcal{R}]^{\mathfrak{n}}$ 的复合（见 (366)）是显然的态射 $\operatorname{Sym}\mathfrak{h}\to[\mathcal{R}/l^{c}(\mathfrak{n})\mathcal{R}]^{\mathfrak{n}}$，其在 $\mathfrak{h}$ 上的限制是 $l^{\mathfrak{b}}$ 与投影 $R_1\to R_1/R_0$ 的复合。

\bdp{7.13.14} \emph{引理.} 设 $l^{c}$ 正则，且态射 $\operatorname{Sym}\mathfrak{h}\to[\mathcal{R}/l^{c}(\mathfrak{n})\mathcal{R}]^{\mathfrak{n}}$ 是同构。则 (370) 是同构。

\emph{证明.} 用 7.13.12(i)。\hfill$\Box$

\bdp{7.13.15} \emph{例.}（参看 [Ko78]）
\begin{enumerate}[label=(\roman*),leftmargin=2.2em]
\item 我们用 7.13.13 的记号。设 $\mathfrak{g}$ 是（有限维）半单李代数，$\mathfrak{b}\subset\mathfrak{g}$ 是 Borel 子代数，$\mathfrak{n}:=[\mathfrak{b},\mathfrak{b}]$。置 $R:=U\mathfrak{g}$，并令 $R_{\cdot}$ 是 $R$ 上的标准滤链，故 $\mathcal{R}=\operatorname{Sym}\mathfrak{g}$。因 $\mathfrak{n}$ 的伴随作用是幂零的，扩张 $\mathfrak{n}^{\flat}$ 典范地平凡化。令 $l:\mathfrak{n}\to\mathfrak{g}\subset R$ 是显然的嵌入。则 $\mathfrak{b}'$ 等于 $\mathfrak{b}\oplus\mathbb{C}$，故此扩张成为平凡的。我们借 7.13.1 的注 (ii) 中的分裂 $s'$ 平凡化扩张 $\mathfrak{b}^{\flat}$。因此我们分裂了扩张 $\mathfrak{b}^{\natural}$，从而 $U^{\natural}\mathfrak{h}=\operatorname{Sym}\mathfrak{h}$。

7.13.14 的条件成立。事实上，$l^{c}$ 显然是正则的，且显然的嵌入 $i^{c}:\operatorname{Sym}\mathfrak{h}\hookrightarrow[\operatorname{Sym}(\mathfrak{g}/\mathfrak{n})]^{\mathfrak{n}}$ 是同构，因为 $\mathfrak{n}$ 沿投影 $(\mathfrak{g}/\mathfrak{n})^{*}\to\mathfrak{h}^{*}$ 的一般纤维单可迁地作用。因此 $h:\operatorname{Sym}\mathfrak{h}\xrightarrow{\sim}H^{0}A$。

下面说明典范态射 (368) $\mathfrak{z}\to H^{0}A=\operatorname{Sym}\mathfrak{h}$ 就是通常的 Harish-Chandra 态射。显然的嵌入 $i:\operatorname{Sym}\mathfrak{h}\xrightarrow{\sim}[R/Rl(\mathfrak{n})]^{\mathfrak{n}}$ 是同构；按定义，Harish-Chandra 态射是嵌入 $\mathfrak{z}\hookrightarrow R^{\mathfrak{n}}$ 与此同构的逆的复合。考虑 (369) 中的映射 $p:H^{0}A\to[R/Rl(\mathfrak{n})]^{\mathfrak{n}}$。由 $p$ 的定义可得 $ph=i$，这蕴含我们的断言。

\SourcePage{344}

\item 现在设 $\psi:\mathfrak{n}\to\mathbb{C}$ 是 $\mathfrak{n}$ 的非退化特征标（我们用 7.13.15(i) 的记号）。置 $R_t:=R[t]$，$l_t:=l+t\psi:\mathfrak{n}\to R_t$。
% [Tutor Note: 原书 p.344 此处作「7.13.15 (i)」，而本条自身即 7.13.15(ii)，所指记号实在 (i) 中给出，编号疑有误（或当为 7.13.13），底本如此，保留原样。]
\end{enumerate}

\bdp{7.13.16} 现在过渡到无穷维情形。我们需要固定一些 Clifford 代数的记号。设 $F$ 是 Tate 向量空间，于是有归纳概形 $Gr(F)$（见 7.11.2(iii)）。归纳概形 $Gr(F)\times Gr(F)$ 带有「相对行列式」的典范线丛 $\lambda$。这是分次线丛，带有典范同构
\begin{equation}\tag{371}
\lambda_{(P,P'')}=\lambda_{(P,P')}\otimes\lambda_{(P',P'')}
\end{equation}
以及对 $P'\subset P$ 的等同 $\lambda_{(P,P')}=\det(P/P')$，它们满足显然的相容性；这里我们假定 $\det(P/P')$ 位于度 $-\dim(P/P')$。

考虑带标准对称形式的 Tate 向量空间 $F\oplus F^{*}$ 与 Clifford 代数 $\mathrm{Cl}=\mathrm{Cl}_F:=\mathrm{Cl}(F\oplus F^{*})$。设 $C$ 是不可约离散 $\mathrm{Cl}$-模\footnote{这里「离散」意指 $C$ 中任一元素的零化子是 $F\oplus F^{*}$ 的开子空间。}。因 $C$ 在张量一个一维向量空间的意义下唯一\footnote{$C$ 同构于费米 Fock 空间 $\varinjlim_U\Lambda^{\cdot}(F/U)\otimes\det(P/U)^{*}$（参看 (182)），其中 $P$ 是 $F$ 中的 c-格，$U$ 取遍 $P$ 的一切 c-子格。}，相应的射影空间 $\mathbb{P}$ 有典范的定义（这是归纳概形）。对任意 c-格 $P\subset F\widehat{\otimes}A$，以 $\lambda^{C}_P$ 记 $C\otimes A$ 中被来自 $P$ 与 $P^{\perp}\subset F^{*}\widehat{\otimes}A$ 的 Clifford 算子零化的元素之集。$A$-子模 $\lambda^{C}_P\subset C\otimes A$ 是一条「线」（即秩 $1$ 的直和项），故 $\lambda^{C}$ 是 $C\otimes\mathcal{O}_{Gr(F)}$ 的线子丛。它定义典范嵌入 $Gr(F)\hookrightarrow\mathbb{P}$。有典范等同
\begin{equation}\tag{372}
\lambda_{(P,P')}=\lambda^{C}_P\otimes(\lambda^{C}_{P'})^{*}
\end{equation}
它与 (371) 相容：若 $P'\subset P$，则同构 $\lambda_{(P,P')}\otimes\lambda^{C}_{P'}\xrightarrow{\sim}\lambda^{C}_P$ 由显然的映射 $\lambda_{(P,P')}=\det(P/P')\to\mathrm{Cl}_F/\mathrm{Cl}_F\cdot P'$ 诱导。

代数 $\mathrm{Cl}$ 带有典范分次，使 $F\subset\mathrm{Cl}^{-1}$、$F^{*}\subset\mathrm{Cl}^{1}$。设 $C^{\cdot}$ 是 $C$ 上与 $\mathrm{Cl}_F$ 的分次相容的分次；它在相差一个平移的意义下\SourcePage{345}%
唯一。于是 $\lambda^{C}$ 是齐次线，且 (372) 是分次线丛的同构。
% [Tutor Note: 原书 p.345 此处印作 "homogenuous"，系 "homogeneous" 之误，底本如此，保留原样。]

\bdp{7.13.17} 以 $\overline{\mathrm{Cl}}^{\cdot}=\overline{\mathrm{Cl}}^{\cdot}_F$ 记 $\mathrm{Cl}^{\cdot}$（作为分次代数）关于由左理想 $\mathrm{Cl}\cdot U$ 生成的拓扑的完备化，其中 $U\subset F\oplus F^{*}$ 是开子空间。于是 $C$ 是离散 $\overline{\mathrm{Cl}}$-模。$\overline{\mathrm{Cl}}$ 的作用给出拓扑分次代数的同构 $\overline{\mathrm{Cl}}^{\cdot}\xrightarrow{\sim}\operatorname{End}^{\cdot}_{\mathbb{C}}C$。

分次代数 $\mathrm{Cl}^{\cdot}$ 有典范滤链 $\mathrm{Cl}^{\cdot}_0=\Lambda^{\cdot}F^{*}\subset\mathrm{Cl}^{\cdot}_1\subset\ldots$（见 7.13.1）。我们把 $\overline{\mathrm{Cl}}^{\cdot}$ 上的滤链 $\overline{\mathrm{Cl}}^{\cdot}_i$ 定义为 $\mathrm{Cl}^{\cdot}_i$ 的闭包。如同 7.13.1，经典 Clifford 代数 $\overline{\mathcal{C}l}^{\cdot}:=\operatorname{gr}\overline{\mathrm{Cl}}^{\cdot}$ 是泊松分次拓扑代数。它带有附加分次 $\overline{\mathcal{C}l}^{\cdot}_i:=\operatorname{gr}_i\overline{\mathrm{Cl}}^{\cdot}$；有
\[
\overline{\mathcal{C}l}^{a}_{i}=\varprojlim_{U,V}\Lambda^{i}(F/U)\otimes\Lambda^{a+i}(F^{*}/V),
\]
其中 $U,V$ 分别是 $F,F^{*}$ 中的 c-格。

以 $E=E_F$ 记 $F$ 的自同态的结合代数。令 $E^{Lie}$ 表示 $E$ 作为李代数。注意 $\overline{\mathcal{C}l}^{0}_1$ 是 $\overline{\mathcal{C}l}^{\cdot}$ 的李子代数，它使 $\overline{\mathcal{C}l}^{-1}_1$ 正规化。$\overline{\mathcal{C}l}^{0}_1$ 在 $\overline{\mathcal{C}l}^{-1}_1=F$ 上的伴随作用把 $\overline{\mathcal{C}l}^{0}_1$ 与 $E^{Lie}$ 等同\footnote{用上面关于 $\overline{\mathcal{C}l}^{0}_1$ 的显式描述。}。置 $E^{\flat}:=\overline{\mathcal{C}l}^{0}_1$；这是 $\mathrm{Cl}^{\cdot}$ 的李子代数，它是 $\overline{\mathcal{C}l}^{0}_1=E^{Lie}$ 被 $\mathbb{C}$ 的中心扩张。

我们看到 $E^{\flat}$ 以与 Clifford 作用相容的方式作用在 $C$ 上；此作用保持 $C$ 上的分次。

接下来几节 7.13.18--7.13.22 给出 $E^{\flat}$ 及其某些子代数的方便描述。读者可以跳过它们，直接看 7.13.23。

\bdp{7.13.18} 这里给出 $E^{Lie}$ 的中心扩张 $E^{\flat}$ 的显式描述，它本质上属于 Tate [T]。

设 $E_+\subset E$ 是由有界算子（即像有界的算子）组成的（双边）理想，$E_-\subset E$ 是由离散算子（即核为开的算子）组成的理想。有 $E_++E_-=E$；置 $E_{tr}:=E_+\cap E_-$。对任意 $A\in E_{tr}$，其迹 $\operatorname{tr}A$ 有定义（若 $U'\subset U\subset F$ 是 c-格，使得
% 尾部接 p.346（本单元在此断处停止）
